\section{Recovery relations and the cost of conditioning}\label{sec:preliminaries}

Write $[m]=\{0,\ldots,m-1\}$. For a matrix of entry width~$b$, identify its raw coordinates with
\[
 V=[d]\times[d]\times[b],\qquad
 \operatorname{addr}(i,j,s)=(id+j)b+s,
 \qquad \Pi(i,j,s)=(j,i,s).
\]
Thus $\Pi$ is an involution, and the raw bit at~$p$ is demanded at physical output address $\operatorname{addr}(\Pi(p))$. In particular, for $p=(i,j,s)$ its raw row is~$i$ and its physical output row is~$j$. These two row coordinates must remain distinct in arguments about actual tape geometry. We sometimes write $\Pi(p)$ for the corresponding output address when the address identification is understood.

\begin{definition}[Storage code]\label{def:pre-storage}
Let $V$ be a finite coordinate set and $N(q)\subseteq V$ a recovery set for every $q\in V$. A storage code over a finite alphabet $\mathcal A$ is a set $\mathcal C\subseteq\mathcal A^V$ such that
\[
 x,y\in\mathcal C,\quad x|_{N(q)}=y|_{N(q)}
 \quad\Longrightarrow\quad x_q=y_q.
\]
Equivalently, each coordinate is a function of its recovery set on~$\mathcal C$. The functions may be nonlinear and may depend on the code. For $a=|\mathcal A|\ge2$, the rate of a nonempty code is $|V|^{-1}\log_a|\mathcal C|$.
\end{definition}

The induced relation on $Q\subseteq V$ is $N_Q(q)=N(q)\cap Q$ for $q\in Q$. A code on this induced relation need not extend to one on~$V$. Conversely, projecting an arbitrary code on~$V$ to~$Q$ need not give a code for $N_Q$: the omitted coordinates must first be fixed.

\begin{lemma}[Paid conditioning]\label{lem:pre-conditioning}
\notformalized
Let $H=V\setminus Q$. Every nonempty fibre of a storage code~$\mathcal C$ obtained by fixing the coordinates in~$H$, and then restricting to~$Q$, is a storage code for~$N_Q$. If every code for~$N_Q$ has at most $a^u$ words, then
\[
 |\mathcal C|\le a^{|H|+u}.
\]
\end{lemma}
\begin{informalproof}
Two words in a fixed fibre agree on~$H$. Agreement on $N(q)\cap Q$ therefore gives agreement on all of~$N(q)$, so it determines~$q$. There are at most $a^{|H|}$ fibres, each of size at most~$a^u$.
\end{informalproof}

\begin{definition}[Paired incidence]\label{def:pre-incidence}
The incidence graph $E(N)$ has vertex set $\{L_p:p\in V\}\sqcup\{R_q:q\in V\}$, with an edge $L_pR_q$ exactly when $p\in N(q)$. The pair $L_p,R_p$ is the \emph{demand pair} for~$p$; it is not an extra graph edge. Let
\[
 b_E(r)=\bigl|\{p\in V:\dist_E(L_p,R_p)\le r\}\bigr|,
\]
where disconnected vertices have distance~$\infty$. A sequence on $|V|\to\infty$ satisfies the \emph{paired-distance condition} if $b_E(r)=o(|V|)$ for every fixed integer~$r$.
\end{definition}

The two incidence degrees encode different restrictions: $\deg(R_q)=|N(q)|$ bounds recovery-set size, whereas $\deg(L_p)$ bounds the reuse of a coordinate across recovery sets. Bounded cone depth immediately gives the former, not the latter. If $N$ is symmetric, the incidence graph is the bipartite double cover of its coordinate graph. In that special case a path from $L_p$ to $R_p$ projects to an odd closed walk based at~$p$. For directed recovery, forgetting which endpoint supplies information loses this interpretation.

\begin{definition}[Index code]
For binary inputs on~$V$, an index code of length~$m$ for~$N$ is a map $G:\{0,1\}^V\to\{0,1\}^m$ such that $G(x)$ together with $x|_{N(q)}$ determines~$x_q$ for every~$q$.
\end{definition}

\begin{lemma}[Fibre bound]\label{lem:pre-fibre}
\notformalized
Every fibre of an index code for~$N$ is a binary storage code for~$N$. Consequently, if all such storage codes have size at most $2^{|V|-s}$, every index code has length at least~$s$.
\end{lemma}
\begin{informalproof}
Inside a fibre the broadcast is fixed, so agreement on the recovery set determines the demanded bit. The at most $2^m$ fibres partition $2^{|V|}$ inputs. Hence $2^{|V|}\le2^m2^{|V|-s}$.
\end{informalproof}

For a native computation, let $I_\Gamma(\ell)$ be the raw input coordinates in the dependency cone of physical output~$\ell$ after a genie~$\Gamma$ is fixed. The relation relevant to transposition is
\begin{equation}\label{eq:transpose-relation}
 N_\Gamma(q)=I_\Gamma(\Pi(q)).
\end{equation}
If the machine correctly transposes every input and the schedule is oblivious, its revealed genie word is an index code for this one relation. If the whole-step value width is~$\beta$, that word has at most $\beta|\Gamma|$ bits. Thus a storage bound of $2^{(1-\rho)L+o(L)}$ contradicts a genie of cost at most~$\eps L$ whenever $\eps<\rho$. The schedule may be studied without any correctness assumption; correctness is used only in this last identification of the demanded output with the raw bit~$q$.

All entropies below are Shannon entropies in bits unless a base is displayed. For a random word supported on a storage code, every recovery relation holds almost surely. Deterministic projection arguments therefore bound arbitrary such distributions as well as uniform distributions on the code. We state cardinality bounds explicitly so that no support-size conclusion is inferred merely from a low entropy under a particular input law.
